% ===========================================================================
% Chapter 3B: Geographic Analysis — Northeast Plain & North China Plain
% ===========================================================================
\section{Study Area: Northeast Plain and North China Plain}

This study targets the two most important grain-producing regions of China:
the \textbf{Northeast China Plain} (东北平原, $43^\circ$--$48^\circ$N) and the
\textbf{North China Plain} (华北平原, $32^\circ$--$40^\circ$N). Together they
account for approximately 45\% of China's total grain output. The contrasting
geographic, climatic, and agricultural characteristics of these two regions
provide an ideal testbed for the mathematical foundations of FusionCropNet V6.

\subsection{Northeast China Plain (东北平原)}

\subsubsection{Geographic Setting}

The Northeast Plain extends across Heilongjiang, Jilin, and Liaoning provinces,
covering approximately 350,000 km$^2$. It is bounded by the Greater Khingan
Range (大兴安岭) to the west, the Lesser Khingan Range (小兴安岭) to the north,
and the Changbai Mountains (长白山) to the southeast. The central plain is
exceptionally flat, with slopes predominantly below $1^\circ$, making it one of
the most mechanized agricultural regions in China.

\begin{table}[h]
\centering
\caption{Northeast Plain: geographic and climatic profile}
\begin{tabular}{lcl}
\toprule
\textbf{Characteristic} & \textbf{Value} & \textbf{Agricultural Significance} \\
\midrule
Latitude & $43^\circ$--$48^\circ$N & Cool-temperate, single cropping \\
Elevation & 50--200m & Minimal terrain correction needed \\
Mean annual temp. & $0^\circ$--$6^\circ$C & Frost-free period 100--160 days \\
Annual precipitation & 400--700mm & Rain-fed agriculture dominant \\
Growing season & May--September & $\sim$150 days, one crop per year \\
Dominant soil & Black soil (Mollisols) & High organic matter (3--6\%), dark spectral signature \\
Field size & 5--50 ha & Large contiguous fields, low edge density \\
Slope & $<1^\circ$ (central plain) & Nearly ideal flat-terrain assumption \\
\bottomrule
\end{tabular}
\end{table}

\subsubsection{Cropping System}

The Northeast Plain practices \textbf{single cropping} (一年一熟) due to the
short frost-free period. The dominant crops are:

\begin{enumerate}[nosep]
    \item \textbf{Spring corn (春玉米):} Planted late April--early May (DOY 110--130),
    harvested late September--early October (DOY 270--290). Accounts for $\sim$40\%
    of China's corn production.
    \item \textbf{Soybean (大豆):} Planted early May (DOY 120--140), harvested late
    September (DOY 260--280). Rotation with corn for nitrogen fixation.
    \item \textbf{Single-cropping rice (单季稻):} Planted mid-May (DOY 130--150),
    harvested late September (DOY 260--280). Concentrated in river valleys with
    irrigation access.
\end{enumerate}

The critical phenological distinction for remote sensing: corn, soybean, and
rice all reach peak NDVI in July--August (DOY 190--240), creating a temporal
classification challenge that the V6 TemporalLite module addresses through
its 3-kernel convolution window covering the critical $\sim$1.5-month
discrimination period.

\subsubsection{Geometric Challenge: Black Soil Spectral Confusion}

The Mollisol (black soil) background presents a unique spectral challenge.
Its high organic matter content (3--6\%) produces reflectance values
significantly lower than the mineral soils of the North China Plain across
all optical bands:

\begin{equation}
    \rho_{\text{black soil}}(\lambda) \approx 0.6 \cdot \rho_{\text{mineral soil}}(\lambda),
    \quad \lambda \in [400, 2500]\text{nm}
\end{equation}

This systematic bias means that an optical classifier trained on the North China
Plain will systematically \textit{underestimate} vegetation vigor in the Northeast
Plain due to darker soil background. The \textbf{DomainAdapter (Module 2)} directly
addresses this cross-region distribution shift: the SGW barycentric mapping
transports source-domain feature distributions (trained on mineral soils) to
target-domain distributions (black soils), without requiring labeled samples
in the target region.

Furthermore, the \textbf{geometric invariants (Module 1)} are invariant to soil
color — Gaussian curvature $K$ and mean curvature $H$ depend only on surface
geometry, not spectral reflectance. This makes DEM-derived features a
\textit{spectrally independent} anchor for cross-region classification, providing
a stable reference when optical features are systematically biased.

\subsubsection{Cloud Cover and SAR Necessity}

The Northeast Plain experiences significant cloud cover during the critical
July--August growing period (60--70\% cloud frequency from MODIS climatology).
This necessitates heavy reliance on Sentinel-1 SAR for temporal continuity.

The \textbf{topological conflict detector (Module 3)} is particularly valuable
here: when optical data is cloud-contaminated, the cohomological conflict
$[\omega_{\text{opt}} \smile \omega_{\text{SAR}}]$ becomes non-trivial,
triggering persistent homology correction that recovers robust consensus
evidence from SAR-dominant observations while appropriately down-weighting
degraded optical evidence.

\subsection{North China Plain (华北平原)}

\subsubsection{Geographic Setting}

The North China Plain extends across Hebei, Shandong, Henan, northern Anhui,
and northern Jiangsu provinces, covering approximately 310,000 km$^2$. It is
an alluvial plain formed by the Yellow River (黄河), Huai River (淮河), and
Hai River (海河) systems. The terrain is predominantly flat ($<1^\circ$ slope)
with local micro-relief variations of 1--5m from ancient river channels and
alluvial fans.

\begin{table}[h]
\centering
\caption{North China Plain: geographic and climatic profile}
\begin{tabular}{lcl}
\toprule
\textbf{Characteristic} & \textbf{Value} & \textbf{Agricultural Significance} \\
\midrule
Latitude & $32^\circ$--$40^\circ$N & Warm-temperate, double cropping \\
Elevation & 0--100m & Coastal to inland gradient \\
Mean annual temp. & $10^\circ$--$15^\circ$C & Frost-free period 190--220 days \\
Annual precipitation & 500--900mm & Strong east-west gradient \\
Growing season & March--November & $\sim$240 days, two crops per year \\
Dominant soil & Alluvial (Inceptisols) & High fertility, variable texture \\
Field size & 1--10 ha & Fragmented landholdings, high edge density \\
Irrigation & Extensive (groundwater) & Over-extraction concerns \\
\bottomrule
\end{tabular}
\end{table}

\subsubsection{Double Cropping System}

The North China Plain practices \textbf{double cropping} (一年两熟), creating
a fundamentally different temporal signature from the Northeast Plain:

\begin{enumerate}[nosep]
    \item \textbf{Winter wheat (冬小麦):} Planted early October (DOY 275--290),
    dormant December--February (vernalization), green-up March (DOY 60--80),
    harvested early June (DOY 150--165). The winter NDVI minimum distinguishes
    wheat from all other crops.

    \item \textbf{Summer corn (夏玉米):} Planted mid-June immediately after wheat
    harvest (DOY 165--175), harvested late September (DOY 265--275). The rapid
    growth phase (V6--VT stages) occurs in July, coinciding with peak rainfall.

    \item \textbf{Cotton (棉花):} Planted mid-April (DOY 100--115), harvested
    October--November (DOY 280--320). Longer growing season, distinct spectral
    signature during boll opening.
\end{enumerate}

\subsubsection{Temporal Challenge: Double-Cropping Temporal Overlap}

The winter wheat--summer corn rotation creates a unique temporal challenge:
in June (DOY 155--175), wheat senescence (decreasing NDVI) temporally overlaps
with corn emergence (increasing NDVI). A single pixel can contain both signals
within a 15-day Sentinel-2 composite, creating spectral mixing that no
single-temporal classifier can resolve.

The \textbf{TemporalLite (Module 4)} directly addresses this: its 3-kernel
depthwise convolution ($k=3$ covering $\sim$45 days) captures the
senescence-to-emergence transition as a single temporal pattern rather than
two conflicting snapshots. The temporal generalization bound
(Eq.~\ref{eq:gen-bound}) guarantees that even with $T_{\text{eff}} < T$
(due to autocorrelation during the stable growing period), the model maintains
bounded generalization error.

\subsubsection{Micro-Relief and Field Boundary Detection}

Although the North China Plain is macroscopically flat, ancient Yellow River
channels create subtle elevation variations of 1--5m over distances of
50--200m. These micro-topographic features serve as natural field boundaries
and drainage divides.

The \textbf{Boundary Detection auxiliary head} in V6 exploits these subtle
DEM gradients: the geometric invariant $\tau_g$ (geodesic torsion) is
particularly sensitive to local surface inflection points, detecting field
boundaries with higher precision than simple Sobel gradient thresholding.
On the North China Plain, where field boundaries often follow micro-relief
rather than visible markers, this geometric boundary detection provides a
critical spatial prior.

\subsection{Cross-Region Comparison and Model Implications}

\begin{table}[h]
\centering
\caption{Contrasting characteristics and V6 module implications}
\begin{tabular}{p{2.8cm}ccp{4.5cm}}
\toprule
\textbf{Characteristic} & \textbf{Northeast} & \textbf{N. China} & \textbf{V6 Module Implication} \\
\midrule
Cropping intensity & Single & Double & \textbf{M4 (TemporalLite):}
$k$=3 captures both patterns; $T_{\text{eff}}$ adjusts automatically \\
Soil background & Dark (Mollisols) & Light (Inceptisols) & \textbf{M2 (DomainAdapter):}
SGW alignment corrects cross-region spectral bias \\
Cloud frequency & High (60--70\%) & Moderate (40--50\%) & \textbf{M3 (Conflict):}
SAR reliance triggers conflict detection; persistent correction recovers consensus \\
Field size & Large (5--50 ha) & Small (1--10 ha) & \textbf{M1 (Boundary):}
$\tau_g$ detects micro-relief boundaries at small fields \\
Terrain variation & Flat ($<1^\circ$) & Flat + micro-relief & \textbf{M1 (Geometric):}
$K\approx0$, but $\tau_g \neq 0$ captures subtle drainage patterns \\
Elevation range & 50--200m & 0--100m & \textbf{M1 (Temporal FiLM):}
Elevation-phenology correction minor but detectable \\
\bottomrule
\end{tabular}
\end{table}

\subsection{Climate-Driven Spectral Dynamics}

\subsubsection{Precipitation Gradient and Vegetation Index Modulation}

The North China Plain exhibits a strong east-west precipitation gradient
(900mm at the coast to 500mm at the piedmont). This gradient directly modulates
NDVI temporal profiles through water availability:

\begin{equation}
    \text{NDVI}_{\max} = \text{NDVI}_{\text{potential}} \cdot
    \left(1 - \exp\left(-\kappa \cdot \frac{P - P_{\min}}{P_{\max} - P_{\min}}\right)\right)
\end{equation}
where $P$ is growing-season precipitation and $\kappa \approx 2.5$ is an
empirical water-use efficiency coefficient calibrated from MODIS NDVI
climatology.

The \textbf{DomainAdapter's running statistics} (Module 2, Section 2.4) track
this precipitation-driven spectral shift during training, enabling the OT
transport plan to adapt continuously rather than requiring separate models for
each precipitation regime.

\subsubsection{Temperature and Phenological Timing}

The latitudinal temperature gradient between the two study regions
($\sim$6$^\circ$C mean annual difference) produces systematic phenological
offsets:

\begin{equation}
    \Delta \text{DOY}_{\text{phenology}} = \frac{\Delta T_{\text{mean}} \cdot
    \text{GDD}_{\text{required}}}{T_{\text{mean}} - T_{\text{base}}}
\end{equation}

For corn (GDD$_{\text{required}} \approx 1500^\circ$C$\cdot$d, $T_{\text{base}} = 10^\circ$C),
a $6^\circ$C mean temperature difference produces approximately a 15--20 day
phenological shift between the two regions. The \textbf{Fourier DOY Encoding}
in TemporalLite captures this shift as a phase offset in the sinusoidal
representation, while the $T_{\text{eff}}$ estimate from Module 4 accounts for
the increased temporal variance.

\subsection{Expected Model Performance by Region}

Based on the geographic analysis above, we formulate testable predictions for
V6 model performance:

\begin{table}[h]
\centering
\caption{Hypothesized relative model performance by region and module}
\begin{tabular}{lcccc}
\toprule
\textbf{Scenario} & \textbf{Baseline} & \textbf{+Module 1} & \textbf{+Module 2} & \textbf{+All Modules} \\
\midrule
Northeast, corn (large fields) & 78\% & 81\% (K discriminates) & 83\% (soil adapted) & \textbf{86\%} \\
Northeast, soybean (dark soil) & 72\% & 76\% ($\tau_g$ boundaries) & 80\% (SGW aligned) & \textbf{83\%} \\
N. China, wheat (winter signal) & 75\% & 77\% (minimal terrain) & 79\% (precip adapted) & \textbf{84\%} \\
N. China, corn (double-crop) & 70\% & 73\% (micro-relief) & 75\% & \textbf{82\%} \\
Cross-region (NE$\to$NC) & 58\% & 65\% (SE(3) invariant) & 72\% (OT transported) & \textbf{78\%} \\
Cross-region (NC$\to$NE) & 55\% & 63\% (SE(3) invariant) & 70\% (OT transported) & \textbf{76\%} \\
\bottomrule
\end{tabular}
\end{table}

The largest expected gains are in cross-region transfer scenarios (NE$\to$NC
and NC$\to$NE), where the SE(3)-invariant geometric features (Module 1) and
SGW domain adaptation (Module 2) provide complementary mechanisms for
overcoming distribution shift. The all-module configuration is predicted to
reduce the cross-region performance gap from $\sim$23\% (baseline) to
$\sim$2--4\%.

\subsection{Remote Sensing Data Specifications}

\subsubsection{Satellite Platforms and Sensors}

\begin{table}[h]
\centering
\caption{Remote sensing data sources}
\begin{tabular}{lcccc}
\toprule
\textbf{Sensor} & \textbf{Platform} & \textbf{Bands} & \textbf{Resolution} & \textbf{Revisit} \\
\midrule
\multicolumn{5}{c}{\textit{Optical}} \\
MSI & Sentinel-2A/B & 13 (VIS+NIR+SWIR) & 10m/20m & 5 days \\
OLI & Landsat 8/9 & 11 (VIS+NIR+SWIR+TIRS) & 30m & 16 days \\
\multicolumn{5}{c}{\textit{SAR (C-band)}} \\
C-SAR & Sentinel-1A/B & VV+VH (IW mode) & 10m & 6--12 days \\
\multicolumn{5}{c}{\textit{DEM}} \\
--- & SRTM v3 & 1 (elevation) & 30m & static \\
--- & ALOS AW3D30 & 1 (elevation) & 30m & static \\
\bottomrule
\end{tabular}
\end{table}

\subsubsection{Spectral Band Utilization}

The 10 optical channels input to V6 are:
\begin{enumerate}[nosep]
    \item Blue (B2, 490nm) — atmospheric correction reference, sensitive to
    aerosol loading differences between the two study regions
    \item Green (B3, 560nm) — vegetation vigor; discriminates corn (deep green)
    from soybean (lighter green) during peak growing season
    \item Red (B4, 665nm) — chlorophyll absorption maximum; winter wheat
    senescence detection in June
    \item Red Edge 1--3 (B5/B6/B7, 705/740/783nm) — critical for crop type
    discrimination at peak biomass when NDVI saturates
    \item NIR (B8, 842nm) — canopy structure; LAI estimation via NDVI
    \item Narrow NIR (B8A, 865nm) — water content; irrigation status in
    the North China Plain
    \item SWIR 1 (B11, 1610nm) — lignin/cellulose; crop residue detection
    after harvest
    \item SWIR 2 (B12, 2190nm) — soil mineral composition; discriminates
    black soil (Northeast) from alluvial soil (North China)
\end{enumerate}

\subsubsection{SAR Backscatter Characteristics}

The 5 SAR input channels are:
\begin{enumerate}[nosep]
    \item VV polarization (dB) — surface scattering; soil moisture after
    rainfall/irrigation events
    \item VH polarization (dB) — volume scattering; canopy structure
    development during vegetative growth
    \item VV/VH ratio — depolarization indicator; discriminates broadleaf
    crops (soybean, cotton) from narrow-leaf crops (wheat, corn)
    \item RVI (Radar Vegetation Index) — $4\sigma_{\text{VH}}/(\sigma_{\text{VV}} +
    \sigma_{\text{VH}})$; robust vegetation proxy in cloudy conditions
    \item Temporal coherence — interferometric stability; detects field
    operations (tillage, harvest) as coherence drops
\end{enumerate}

In the Northeast Plain, SAR is essential during the July--August cloudy period
when optical data is unreliable. In the North China Plain, SAR provides
critical observations during the June wheat harvest--corn planting transition,
when rapid field changes occur within a single Sentinel-2 revisit interval.
% (end of SAR characteristics section)
